\section{Nonvanishing along the prime sequence}
\label{sec:nonvanishing}

The next proposition supplies the nonzero determinants needed in
Proposition~\ref{prop:finite-lower}.

\begin{proposition}\label{prop:nonvanishing}
Assume that $G\in\mathbb Q$. For every sufficiently large prime $p$, the
determinant with scale parameter $N=p$ satisfies
\begin{equation}\label{eq:nonvanishing-exact-valuation}
 v_p(\Delta_p)=-96p.
\end{equation}
In particular, $\Delta_p\ne0$ for every sufficiently large prime $p$.
\end{proposition}

We prove the proposition by reducing a matrix of size $48p$ to three fixed
rational matrices of size $48$. Throughout this section, an underlined row
polynomial is formed with $N=1$, so that
\[
 n_0=48,\quad b_0=7,\quad q_0=g_0=4,\quad h_0=2,\quad
 L_0=59,\quad C_0=63,\quad H_0=65.
\]
In addition to the usual rows $0,\ldots,47$, define the auxiliary row $48$
by exactly the formulas in Equation~\eqref{eq:rows}. For $0\le r\le48$ and
$0\le k<48$, put
\begin{equation}\label{eq:nonvanishing-base-matrix}
 \mathcal B(r,k)=\sum_{j=0}^{4}(-1)^j\binom4j
 \left\{M^0(\underline P_r,7+k+j)
              -\frac32 Z^0(\underline D_r,7+k+j)\right\}.
\end{equation}
Here $M^0$ and $Z^0$ act linearly on the first polynomial argument. Thus
$\mathcal B$ is a fixed rational $49$ by $48$ matrix. Define
\[
 \mathcal B_0=(\mathcal B(r,k))_{0\le r,k<48},\qquad
 \mathcal B_1=(\mathcal B(r+1,k))_{0\le r,k<48}.
\]
The arithmetic certificate below proves
\begin{equation}\label{eq:nonvanishing-fixed-determinants}
 \det\mathcal B_0\ne0,\qquad
 \det(\mathcal B_0+\mathcal B_1)\ne0,\qquad
 \det(\mathcal B_0-\mathcal B_1)\ne0.
\end{equation}
First we show why these three fixed assertions imply the proposition.

\subsection{Two palindromic bases}

Let $p$ be an odd prime and work over $\mathbb F_p$. The vector space
\[
 \mathcal P_p=\{Q\in\mathbb F_p[w]:\deg Q\le2p-2,
                \ w^{2p-2}Q(1/w)=Q(w)\}
\]
has dimension $p$: its coefficients of $w^0,\ldots,w^{p-1}$ determine
all its coefficients. For $0\le\ell,i<p$, consider
\[
 U_\ell(w)=(2w)^\ell(1+w^2)^{p-1-\ell},\qquad
 Q_i(w)=w^i\sum_{j=0}^{p-i-1}w^{2j}.
\]
Both families belong to $\mathcal P_p$. The lowest terms of $U_\ell$
and $Q_i$ are respectively $2^\ell w^\ell$ and $w^i$. Triangularity
of their coefficients in degrees $0,\ldots,p-1$ proves that each family
is a basis. Define $a_{\ell i}\in\mathbb F_p$ by
\[
 Q_i=\sum_{\ell=0}^{p-1}a_{\ell i}U_\ell,
 \qquad \mathbf a=(a_{\ell i})_{0\le\ell,i<p}.
\]
The same lowest-term comparison gives
\begin{equation}\label{eq:nonvanishing-triangular-unit}
 a_{\ell i}=0\quad(\ell<i),\qquad a_{ii}=2^{-i},\qquad
 \det\mathbf a=2^{-p(p-1)/2}\ne0.
\end{equation}
In particular, invertibility of this varying-size matrix introduces no
exceptional odd primes.

Define $Q'_0=0$ and $Q'_i=Q_{p-i}$ for $1\le i<p$, and write
$Q'_i=\sum_\ell b_{\ell i}U_\ell$. Equivalently,
\begin{equation}\label{eq:nonvanishing-transition}
 \mathbf b=\mathbf a\Pi,\qquad
 \Pi e_0=0,\qquad \Pi e_i=e_{p-i}\quad(1\le i<p),
\end{equation}
where $e_i$ are the standard coordinate vectors. The explicit forms are
\[
 Q_i=w^i\frac{1-w^{2(p-i)}}{1-w^2},\qquad
 Q'_i=w^{p-i}\frac{1-w^{2i}}{1-w^2};
\]
the latter formula gives zero also when $i=0$. Their sum satisfies
\[
 Q_i+w^pQ'_i
 =w^i\sum_{j=0}^{p-1}w^{2j}
 =w^i(1-w^2)^{p-1},
\]
because $\binom{p-1}{j}=(-1)^j$ in $\mathbb F_p$. With
\[
 t=\frac{2w}{1+w^2},\qquad f=\frac{1-w^2}{1+w^2},\qquad
 E(t)=(1-t^2)^{(p-1)/2}=f^{p-1},
\]
division by $(1+w^2)^{p-1}$ therefore proves the rational-function identity
\begin{equation}\label{eq:nonvanishing-palindromic-expansion}
 E(t)w^i=\sum_{\ell=0}^{p-1}t^\ell
              (a_{\ell i}+b_{\ell i}w^p).
\end{equation}

\subsection{Frobenius and the two extraction shifts}

Set $N=p$, and write every row index uniquely as
$r=pr_0+i$, with $0\le r_0<48$ and $0\le i<p$. Put
$t'=t^p$, $w'=w^p$, and $f'=f^p$. Frobenius gives
\[
 t'=\frac{2w'}{1+(w')^2},\qquad
 f'=\frac{1-(w')^2}{1+(w')^2},\qquad (f')^2=1-(t')^2.
\]
Using $C=63p$, $h=2p$, and $g=4p$ in the row definition gives
\begin{align}
 tR_rE(t)
 &= (1-t')^2(t')^{63}(w')^{r_0-4}E(t)w^i\notag\\
 &=\sum_{\ell=0}^{p-1}t^\ell t'
       \bigl(a_{\ell i}\underline R_{r_0}(t',w')
            +b_{\ell i}\underline R_{r_0+1}(t',w')\bigr).
 \label{eq:nonvanishing-frobenius-row}
\end{align}
Here the single factor $w'$ in the second term creates exactly one successor
row; its index is at most $48$.

To separate the polynomial parts, star fixes $t,t'$ and negates $f,f'$.
Moreover $fE=f'$, so the row identities give
\[
 tR_rE=tP_rE-f'D_r,
 \qquad t'\underline R_j=t'\underline P_j-f'\underline D_j.
\]
Taking the star-invariant part of
Equation~\eqref{eq:nonvanishing-frobenius-row} and then its anti-invariant
part, and cancelling the nonzero rational function $f'$ in the latter,
yields two separate polynomial identities:
\begin{align}
 tP_r(t)E(t)
 &=\sum_{\ell=0}^{p-1}t^\ell t'
       \bigl(a_{\ell i}\underline P_{r_0}(t')
              +b_{\ell i}\underline P_{r_0+1}(t')\bigr),
       \label{eq:nonvanishing-P-identity}\\
 D_r(t)
 &=\sum_{\ell=0}^{p-1}t^\ell
       \bigl(a_{\ell i}\underline D_{r_0}(t')
              +b_{\ell i}\underline D_{r_0+1}(t')\bigr).
       \label{eq:nonvanishing-D-identity}
\end{align}
These are polynomial identities since the row polynomials on both sides
are polynomials, and substitution $t=2w/(1+w^2)$ is injective on
$\mathbb F_p[t]$.

In the notation of Equation~\eqref{eq:column-layer}, the extractions for a
fixed residue $\ell$ are
\[
 P'_u=[t^{(u+1)p+\ell}]tP_r(t)E(t),\qquad
 D'_u=[t^{up+\ell}]D_r(t).
\]
Every polynomial has a unique expression
$\sum_{\ell=0}^{p-1}t^\ell A_\ell(t^p)$. Consequently
Equations~\eqref{eq:nonvanishing-P-identity} and
\eqref{eq:nonvanishing-D-identity} give
\begin{equation}\label{eq:nonvanishing-extracted-rows}
 (P'_u,D'_u)=[(t')^u]
 \left\{a_{\ell i}(\underline P_{r_0},\underline D_{r_0})(t')
        +b_{\ell i}(\underline P_{r_0+1},\underline D_{r_0+1})(t')\right\}.
\end{equation}
Thus the $t'$ in the first identity is exactly consumed by the $u+1$
in the $P$ extraction. The $D$ extraction has no such shift. Negative
coefficient indices are zero; in particular $P'_{-1}=0$ here.

There is a degree issue at the last successor row which is useful to make
explicit. The auxiliary polynomials have
\begin{equation}\label{eq:nonvanishing-auxiliary-degrees}
 \min\deg\underline P_{48}=62-44=18,\qquad
 \min\deg\underline D_{48}=63-44=19.
\end{equation}
These coefficients must be retained, even though the first $48$ base rows
have minimum degrees at least $19$ and $20$. For the last block
$r=47p+i$, $i>0$, Equation~\eqref{eq:nonvanishing-transition} gives
$b_{\ell i}=a_{\ell,p-i}=0$ for $\ell<p-i$, and
$b_{p-i,i}=2^{-(p-i)}\ne0$. The successor contribution in either
$tP_rE$ or $D_r$ therefore starts at
\[
 19p+(p-i)=20p-i.
\]
The first-row contribution starts no earlier than $20p+i$.
Directly from the original rows, $|r-g|=43p+i$, so $tP_r$ and $D_r$
also start at $C-|r-g|=20p-i$; since $E(0)=1$, this agrees with the
extraction. The leading Chebyshev coefficients are powers of $2$, hence
are nonzero in odd characteristic. When $i=0$, the successor term is
zero and the first term starts at $20p$. The lower degree of the
auxiliary row therefore introduces no forbidden coefficient into the
growing matrix. Its contact identity is also valid:
$t\underline R_{48}/f=O(t^{107})$, since $63+48-4=107>59$.

\subsection{The residue blocks and their determinant}

For this paragraph take $p>260$ and exclude primes dividing the denominator
of the hypothesized rational number $G$. Then
\[
 p>2\sqrt{65p}=2\sqrt H,
\]
so Equations~\eqref{eq:digit-layers} and \eqref{eq:column-layer} apply at
the same prime $p$ used as the scale parameter. To spell out the role of
$G$, the rational raw entry is
\[
 F_j(r)=M^0(P_r,j)-\frac32 Z^0(D_r,j)
          +4G\sum_v[t^v]P_r\,c_{(v-j)/2}.
\]
The last sum is $p$-integral: its polynomial coefficients are integers,
and every $c_d=4^{-d}\binom{2d}{d}$ is integral at an odd prime.
The $G$ term therefore reduces to zero after multiplication by $p^2$.
All entries of the scaled matrix are $p$-integral by the cited layer
formulas. Substitution of Equation~\eqref{eq:nonvanishing-extracted-rows}
into Equation~\eqref{eq:column-layer} gives the reduced raw entry
\begin{align}
 p^2F_{kp+\ell}(pr_0+i)
 \equiv{}&a_{\ell i}\left\{M^0(\underline P_{r_0},k)
                         -\frac32 Z^0(\underline D_{r_0},k)\right\}
                         \notag\\
 &+b_{\ell i}\left\{M^0(\underline P_{r_0+1},k)
                         -\frac32 Z^0(\underline D_{r_0+1},k)\right\}
                         \pmod p.
 \label{eq:nonvanishing-raw-bridge}
\end{align}
This applies to every raw column in the specified range. The base
entries here have denominators with prime factors at most $65$, as
the rational recipes below also show, so their reductions exist for
$p>260$.

The integer filter also respects residues. Write the column index
uniquely as $k=\ell+pk_0$, where
$0\le\ell<p$ and $0\le k_0<48$. In $\mathbb F_p[s]$,
\begin{equation}\label{eq:nonvanishing-residue-filter}
 s^{7p+k}(1-s)^{4p}
   =s^\ell(s^p)^{7+k_0}(1-s^p)^4.
\end{equation}
Multiplying each raw entry by $p^2$ before reduction makes this polynomial
congruence applicable to the column operation: coefficient differences
divisible by $p$ multiply integral scaled entries. The fourth finite
difference on the right uses exactly the five base raw indices
$7+k_0,\ldots,11+k_0$,
all in $7,\ldots,58$. Each of the $p$ residue groups thus contains
exactly $48$ base filtered columns, and no operation changes $\ell$.

Order rows by $(i,r_0)$ and columns by $(\ell,k_0)$. The reduction of
the scaled $48p$ by $48p$ matrix has $(i,\ell)$ block
\[
 a_{\ell i}\mathcal B_0+b_{\ell i}\mathcal B_1.
\]
Writing this matrix as $\mathcal L_p$ makes its orientation explicit:
\begin{equation}\label{eq:nonvanishing-block-matrix}
 \begin{split}
 \mathcal L_p&=\mathbf a^T\otimes\mathcal B_0
                       +\mathbf b^T\otimes\mathcal B_1,\\
 \mathcal L_p\bigl((\mathbf a^T)^{-1}\otimes I_{48}\bigr)
      &=I_p\otimes\mathcal B_0+\Pi^T\otimes\mathcal B_1.
 \end{split}
\end{equation}
Indeed $\mathbf b^T=\Pi^T\mathbf a^T$, which explains both the
transpose and the side of multiplication.

The nonzero indices $1,\ldots,p-1$ fall into $(p-1)/2$ disjoint pairs
$\{i,p-i\}$. On each pair the vectors $e_i+e_{p-i}$ and
$e_i-e_{p-i}$ are eigenvectors of $\Pi$ with eigenvalues $1$ and
$-1$. Together with $e_0$, of eigenvalue zero, they form a basis since
$2$ is invertible. Thus the eigenvalue multiplicities of $\Pi^T$ are
$1,(p-1)/2,(p-1)/2$ for $0,1,-1$, respectively. A similarity on the
$p$-dimensional factor in Equation~\eqref{eq:nonvanishing-block-matrix}
now proves the exact identity in this block ordering:
\begin{align}
 \det\mathcal L_p
  ={}&(\det\mathbf a)^{48}\det\mathcal B_0\notag\\
    &\quad\cdot\det(\mathcal B_0+\mathcal B_1)^{(p-1)/2}
          \det(\mathcal B_0-\mathcal B_1)^{(p-1)/2}.
 \label{eq:nonvanishing-factorization}
\end{align}
There is no determinant factor from the similarity. Returning to the
original row and column orders changes at most the sign.

\subsection{An exact certificate for the fixed matrices}

All entries in Equation~\eqref{eq:nonvanishing-base-matrix} can be
constructed by the following rational arithmetic. This also specifies
their reduction modulo $101$ without any numerical approximation.
First form $m_0=2$, $m_i=0$ for odd $i$, and
$m_i=i m_{i-2}/(i+1)$ for even $2\le i\le64$. Set
\[
 k^-_0=0,\quad k^-_1=2,\qquad k^+_0=k^+_1=0,
\]
and use
\begin{align*}
 k^-_d&=\frac{(d-1)k^-_{d-2}+m_{d-2}+m_{d-1}}{d}
                  &&(2\le d\le64),\\
 k^+_d&=\frac{(d-2)k^+_{d-2}+2/(d-1)}{d-1}
                  &&(2\le d\le58).
\end{align*}
On $0\le i\le64$, $0\le j\le58$, form
\[
 Y_{i0}=k^-_i,\qquad Y_{0j}=k^+_j,\qquad
 Y_{ij}=Y_{i-1,j-1}-\frac{m_{i-1}}j\quad(i,j\ge1).
\]
Thus $Y_{ij}=M^0(i,j)$. With $B_i^{(d)}=\sum_{u=1}^i u^{-d}$ and
$B_0^{(d)}=0$, set
\[
 J_{ij}=\frac32\begin{cases}
       -B_i^{(2)},&i=j,\\[2pt]
       (B_i^{(1)}-B_j^{(1)})/(i-j),&i\ne j.
      \end{cases}
\]
This gives $J_{ij}=\tfrac32Z^0(i,j)$, with the sign on its diagonal
included.

For the polynomial construction define
\[
 E_0=t^{62},\quad E_1=t^{61},\qquad I_0=0,\quad I_1=t^{62},
\]
and, for either family $S=E,I$, use
\[
 S_m=2S_{m-1}/t-S_{m-2}\qquad(2\le m\le44).
\]
These are ordinary integer polynomials: they are
$E_m=t^{62}T_m(1/t)$ and $I_m=t^{62}U_{m-1}(1/t)$.
For $0\le r\le48$ let $u=|r-4|$ and form the coefficient vectors,
in degrees $0,\ldots,64$, of
\[
 y_r(t)=(1-t)^2E_u(t),\qquad
 z_r(t)=\operatorname{sign}(r-4)(1-t)^2I_u(t).
\]
They are exactly $\underline P_r$ and $\underline D_r$.
Contract them with the arrays to obtain the length-$52$ vector
\[
 v_r(j)=\sum_{i=0}^{64}\bigl([t^i]y_r\,Y_{ij}
                                  -[t^i]z_r\,J_{ij}\bigr),
               \qquad j=7,\ldots,58.
\]
Replace a vector $v$ four times successively by its vector of consecutive
differences $(v(j)-v(j+1))_j$. The resulting length-$48$ vector is
row $r$ of $\mathcal B$. In particular, the construction includes all
coefficients of the auxiliary row, including its degree-$18$ term.

Every scalar denominator in these recipes is a product of nonzero integers
of absolute value at most $65$. Every rational array entry therefore has
denominator prime factors at most $65$, so each denominator is a unit
modulo $101$. Polynomial division by $t$ above shifts exponents of
polynomials divisible by $t$ and introduces no scalar denominator.
Consequently the entire construction can be performed over
$\mathbb F_{101}$ and agrees with reduction of the rational matrices.

For completeness, the elimination rule producing the certificate is as
follows. For $\sigma\in\{0,1,-1\}$, start with row vectors
$R_i=\mathcal B(i,{\cdot})+\sigma\mathcal B(i+1,{\cdot})$,
$0\le i<48$, in increasing order. At step $i$, if $R_i(i)=0$, swap
row $i$ with the first later row having a nonzero entry in column $i$.
Record $d_i=R_i(i)$ and replace every row $j>i$ by
\[
 R_j-\frac{R_j(i)}{d_i}R_i.
\]
There are no swaps for $\sigma=0,-1$. For $\sigma=1$ the only swaps,
using indices starting at zero, are $(30,31)$ and $(45,46)$ at their
respective steps. The pivots are given in
Table~\ref{tab:modular-pivots}; its three lines for each $\sigma$
list the pivots consecutively, with $16$ entries per line.

\begin{table}[htbp]
\centering
\small
\begin{tabular}{cl}
\toprule
$\sigma$ & Successive pivots in $\mathbb F_{101}$\\
\midrule
$0$ & $60,68,62,79,47,32,69,57,30,72,35,66,43,20,85,48$\\
& $88,4,77,54,60,79,26,68,83,39,40,65,1,68,78,24$\\
& $15,98,32,22,94,9,99,10,15,75,4,2,25,53,90,79$\\[2pt]
$1$ & $38,51,90,70,5,21,88,55,45,20,35,41,77,10,18,25$\\
& $76,14,38,72,6,66,56,35,83,58,56,11,17,20,30,24$\\
& $28,11,25,70,79,99,66,38,4,41,63,91,63,17,90,98$\\[2pt]
$-1$ & $82,92,26,87,21,74,87,88,88,3,14,23,38,58,36,20$\\
& $26,33,94,74,78,45,93,86,73,66,45,30,61,3,88,27$\\
& $20,58,69,48,78,39,48,1,66,18,86,93,52,92,39,77$\\
\bottomrule
\end{tabular}
\caption{Exact Gaussian pivots for $\mathcal B_0+\sigma\mathcal B_1$
modulo $101$.}\label{tab:modular-pivots}
\end{table}

Every listed pivot is nonzero. Since the swaps and row additions are
invertible, all three matrices are invertible modulo $101$ and hence
have nonzero determinants over $\mathbb Q$. This proves
Equation~\eqref{eq:nonvanishing-fixed-determinants}. The modulus $101$
is used only for this fixed finite certificate; the growing determinant
uses arbitrary sufficiently large primes as follows.

\begin{proof}[Completion of the proof of Proposition~\ref{prop:nonvanishing}]
Let $\mathcal E$ consist of $2$, the primes dividing the denominator of
$G$, the primes dividing any denominator of an entry of $\mathcal B$,
and the primes dividing the numerator or denominator, in lowest terms,
of any of the three nonzero rational determinants in
Equation~\eqref{eq:nonvanishing-fixed-determinants}. This is a fixed
finite set. For $p>260$ outside $\mathcal E$, all three fixed matrices
reduce to invertible matrices over $\mathbb F_p$. Equation~
\eqref{eq:nonvanishing-triangular-unit} and the factorization in
Equation~\eqref{eq:nonvanishing-factorization} then show that
$\det\mathcal L_p\ne0$.

Let $\mathcal F_p$ be the original filtered matrix whose determinant is
$\Delta_p$. Its size is $n=48p$, and $p^2\mathcal F_p$ is integral
over the local ring $\mathbb Z_{(p)}$. Its reduction, after the stated
row and column permutations, is $\mathcal L_p$. Therefore
\[
 p^{2n}\Delta_p=\det(p^2\mathcal F_p)
\]
is a unit in $\mathbb Z_{(p)}$. This proves
$v_p(\Delta_p)=-2n=-96p$. Every sufficiently large prime is outside
$\mathcal E$ and exceeds $260$, as required.
\end{proof}
